\begin{resumo}[Abstract]
 \begin{otherlanguage*}{english}
   Project EMA, developed at the Automation and Robotics Laboratory (LARA) of the University
   of Brasília, builds motor rehabilitation systems that integrate functional electrical
   stimulation, robotics and virtual reality. Over the years, successive generations
   of students have produced several subsystems that, although functional from a research
   standpoint, exhibit recurring software engineering weaknesses: coupled architectures,
   lack of documentation, disorganized version control and a systematic gap between what is
   recorded and what is actually assembled. This applied, exploratory work applies software
   engineering methods to the diagnosis, documentation and restructuring proposal of these
   systems, adopting a consulting stance: it takes the perspective of a new member who tries
   to reproduce the systems as they are, recording every obstacle as diagnostic data. The
   investigation is organized around two complementary case studies at opposite maturity
   stages: the immersive virtual reality system, inactive for years, whose
   challenge is reactivation; and the functional electrical stimulation cycling system,
   active yet hard to reproduce. For each case, the current architecture is documented, the
   reproduction attempt is reported and the specific limitations are mapped; the common
   problems are then consolidated, with emphasis on version control organization and
   documentation. As a result, the work delivers a structured diagnosis, standardized
   architecture diagrams and a set of restructuring recommendations, preparing the ground
   for the integration stage to be carried out in future work.

   \vspace{\onelineskip}

   \noindent
   \textbf{Key-words}: software engineering. rehabilitation robotics. documentation. refactoring. Project EMA.
 \end{otherlanguage*}
\end{resumo}
